
As discussed in the Introduction, our work targets an active learning
setup where the system is trained and deployed and where, during
deployment, the system may request to have a limited number of new
instances manually annotated for the purpose of fine-tuning and
adapting.

We will establish the following nomenclature, reflecting how we mapped
the general deployment scenario above into a concrete ML experiment:
%
\begin{itemize}
\item The \term{training set} and the \term{validation set}
  constitute the train/test split of the original data.
  The training set is used to train a classifier and the
  validation set is used to estimate model accuracy and HAS scale.
  This step reflects model preparation before deployment.
\item The \term{fine-tuning set} and the \term{test set} constitute
  the train/test split of a second dataset. There is an
  \quoted{annotation budget} for labelling some of the instances of
  the fine-tuning set, the \term{labelled fine-tuning set}.
  The test set is fully labelled.
  The success of the selection is estimated by measuring accuracy on
  the test set: In the presence of drift, we expect the accuracy of
  the original model to drop and this drop to be, at least partially,
  recovered by the fine tuned model.
  This step reflects model adaptation after deployment.
  % The \term{fine-tuning set} and the \term{test set}
  % constitute the train/test split of a second, target dataset
  % that represents the post-deployment distribution.
  % The fine-tuning set is first passed through the drift
  % detection pipeline, which assigns each sample a drift
  % category (\emph{In-Distribution}, \emph{Pure Data Drift},
  % \emph{Pure Concept Drift}, or \emph{Full Drift}) based on
  % the geometric signals described in Section~\ref{sec:detection}.
  % Given an \term{annotation budget} \textit{i.e.} a constraint on how
  % many fine-tuning samples may be labelled, the
  % \term{labelled fine-tuning set} is constructed by selecting
  % and prioritising samples according to their drift type and
  % severity, rather than at random. The success of the selection is estimated by measuring accuracy on
  %the test set: In the presence of drift, we expect the accuracy of
  %the original model to drop and this drop to be, at least partially,
  %recovered by the fine tuned model.
\end{itemize}



Since we have scoped our work as addressing gradual drift in the
meaning of existing labels (as opposed to novelty and to abrupt
drift), we assume that the labelled fine-tuning set should
include the instances that are further away from the semantic core of
our labels, since the instances at the periphery are the ones that
establish the decision boundaries.

In the HASeparator embedding space, peripherality is naturally
quantified by the angular margin $\Delta_i$
(Equation~\ref{eq:margin}).
Instances with small $\Delta_i$ lie close to the separation hyperplane 
between their predicted class and the class with the closest decision boundary,
 and are therefore the most
geometrically fragile under gradual drift. The training margin
distribution $\{\Delta_i^{\mathrm{train}}\}$ defines what
constitutes a central versus peripheral instance on
$\mathcal{S}^{63}$
Central instances have large margins,
deep within their class region, while peripheral instances
have small margins, near the decision boundaries.

In this context we formulate the following hypotheses:
%
\begin{itemize}
\item[\textbf{H1}] We will use HAS space and HAS scale to define a
  distribution curve upon which we can place instances to measure how
  central or peripheral they are. This is analogous to how data drift
  is estimated, but we hypothesise that: The instances that are at the
  edges of the the HAS curve more likely to be misclassified even when
  they are \emph{not} at the edges of the data drift curve.

\end{itemize}

We consider that such instances are the primary drivers of
concept drift and that, if H1 holds, they constitute the
most informative subset for fine-tuning. We therefore
further hypothesise:

\begin{itemize}
\item[\textbf{H2}] If we use the HAS space and HAS scale to find the
  peripheral instances, we get a better selection than what we get in
  the original space.
  % Selecting fine-tuning instances by
  % prioritising those with the smallest angular margin
  % \textit{i.e.}, those closest to decision boundaries
  % in HAS space, yields greater accuracy and
  % margin improvement after fine-tuning than an equivalent
  % budget of randomly selected instances or instances selected
  % by confidence score in the original embedding space.
\end{itemize}

%
\begin{itemize}
\item classifier, original space
\item HASeparator, embedding
\item extract: some statistics about the HAS space
\end{itemize}

deploy

H1 is evaluated through the drift taxonomy and population-level
statistical tests described in Section~\ref{sec:detection},
which assess whether margin-peripheral instances exhibit
disproportionate misclassification rates independently of
their data drift status.
H2 is evaluated through the
two-pool fine-tuning experiment of Section~\ref{sec:finetune},
which directly compares the targeted concept-drift selection
strategy (Pool B) against random global selection and
data-drift-only selection (Pool A) under identical annotation
budgets.



test:
%
\begin{itemize}
\item for each instance x class, calculate a valuation based on extract
\item if two classes are near, this is a candidate for represent drift
\item assign a label
\item fine tuning
\end{itemize}





\subsection{Feature Extraction and Geometric Signal Computation}
\label{sec:extraction}


\subsubsection{Baseline Reference Distribution}

For the baseline, two statistics are computed over the training
embeddings and stored as the reference distribution.
The global centroid $\bm{c} = \frac{1}{n}\sum_i \be_i$
and the associated Euclidean distance
$d_i = \|\be_i - \bm{c}\|_2$
characterise the spatial spread of the training latent space.
Prediction confidence $q_i = \max_j\, p_{ij}$, where
$p_{ij}$ is the softmax probability of class $j$, captures
output certainty.
The per-split mean and standard deviation of both $d$ and $q$
serve as calibration statistics for thresholding at test time.

\subsubsection{HAS-Specific Drift Signals}

The spherical geometry of HAS embeddings enables three additional
drift-sensitive signals that are not presented in the baseline setting due to its geometrical limitations. These signals arise directly from the hyperspherical structure that the
HASeparator penalty imposes during training. 
By explicitly pushing embeddings away from class separation hyperplanes, the trained
model encodes proximity to decision boundaries as a measurable
geometric quantity that can be monitored at inference time.
Crucially, these signals are entirely independent of the
HASeparator training hyperparameters ($m$ and $\sigma$), which
govern the geometry during learning.
Thus the drift signals are computed post-hoc from the frozen model's outputs alone.


\paragraph{Angular Margin}

Let $\hat{y}_i$ and $\hat{y}_i^{(2)}$ denote the predicted
(highest-scoring) and second-best (runner-up) class for sample
$i$, respectively, where the runner-up class is the one whose
weight vector $\hat{\bw}_{\hat{y}_i^{(2)}}$ has the second-highest
cosine similarity with the embedding \textit {i.e.}, the class the model
was most tempted to predict instead of $\hat{y}_i$.
The angular margin is defined as the difference between these two
cosine similarities:
\begin{equation}
  \Delta_i \;=\;
    \bigl(\hat{\be}_i \cdot \hat{\bw}_{\hat{y}_i}\bigr) -
    \bigl(\hat{\be}_i \cdot \hat{\bw}_{\hat{y}_i^{(2)}}\bigr)\,,
  \label{eq:margin}
\end{equation}
and corresponds to the angular gap between the embedding and the
nearest class separation hyperplane on $\mathcal{S}^{63}$.
A large $\Delta_i$ indicates that the embedding lies deep within
the predicted class region, comfortably distant from any decision
boundary.
A small positive $\Delta_i$ signals that the embedding is close to the boundary between the predicted and runner-up class, and therefore geometrically fragile under distribution
shift.
Moreover, $\Delta_i \approx 0$ means the embedding sits almost
exactly on the separation hyperplane, whilea negative $\Delta_i$
indicates that the embedding has crossed the boundary, ~\textit{ i.e.}, the
runner-up class is actually more aligned with the embedding than
the predicted class, which is a strong indicator of
misclassification.

Since the HASeparator penalty during training explicitly maximises
this margin for all training samples, the training distribution
$\{\Delta_i^{\mathrm{train}}\}$ provides a natural reference for
what constitutes a geometrically safe embedding.
A custom-set sample is flagged as concept-drifted when its margin falls below:
\begin{equation}
  \tau_\Delta \;=\; \max\!\bigl(\mu_\Delta - k\,\sigma_\Delta,\ 0\bigr),
  \label{eq:margin_thresh}
\end{equation}
where $\mu_\Delta$ and $\sigma_\Delta$ are the mean and standard
deviation of the training margin distribution, and $k = 2.0$
is a sensitivity hyperparameter. A custom sample falling below
$\tau_\Delta$ is therefore one whose embedding is more
boundary-proximate than approximately $97.7\%$ of training
samples, constituting a statistically motivated criterion for
geometric concept drift.



\paragraph{Cosine Distance from the Per-Class Sphere Centroid.}


While the angular margin captures proximity to decision boundaries,
a complementary signal is needed to detect shifts in the input
representation itself ~\textit{i.e.} data-drift.
For this purpose, the cosine distance from the
predicted-class sphere centroid $\bm{\mu}_{\hat{y}_i}$ is used
as the data drift indicator. For a test sample $i$ predicted as
class $\hat{y}_i$, this is defined as:
\begin{equation}
  d_{\mu}(i) \;=\; 1 - \hat{\bm{e}}_i \cdot \bm{\mu}_{\hat{y}_i}\,,
  \label{eq:cos_dist}
\end{equation}
where $\bm{\mu}_{\hat{y}_i}$ is the per-class sphere centroid
computed from all training embeddings of class $\hat{y}_i$.
A low $d_\mu(i)$ indicates that the embedding is tightly
clustered with the training samples of its predicted class on
$\mathcal{S}^{63}$.
Additionally, a high value indicates that the embedding
has drifted away from where training samples of that class
typically reside, even if the model still predicts that class
with high confidence.

This method has two properties that make it particularly suited
to the HAS geometry. First, it is class-conditional.
Thus the reference centroid $\bm{\mu}_{\hat{y}_i}$ depends on the
predicted class, so the threshold adapts to the natural spread
of each class on the hypersphere rather than applying a single
global criterion. Second, it is scale-free, because both
the embedding and the centroid are unit vectors, the cosine
distance is bounded in $[0, 2]$ and is not affected by the
feature scaler $\sigma$ used during training. 



\paragraph{Closest-Boundary Class.}
\label{sec:closest_boundary}
Beyond detecting \emph{whether} a sample is concept-drifted,
the HAS geometry also reveals \emph{toward which class} the
embedding is drifting. For each sample $i$, the closest-boundary
class index (already computed during the
angular margin calculation) is recorded as the
\emph{closest-boundary class}, or equivalently the predicted
drift destination.
This identifies the class whose weight vector is most
aligned with the embedding after the predicted class, and
therefore indicates the decision boundary that the embedding
is closest to on $\mathcal{S}^{63}$.

The closest-boundary class provides two complementary pieces
of information. First, at the \emph{sample level}, it
identifies which specific boundary a concept-drifted sample
is pointing at, enabling interpretable per-instance
drift diagnosis.
Second, at the \emph{population level}, aggregating
$\hat{y}_i^{(2)}$ across all concept-drifted samples
yields a drift direction matrix that reveals systematic
confusion patterns between class pairs under distribution
shift. 


\subsection{Drift Detection}
\label{sec:detection}

Drift detection proceeds at two complementary levels.
A population-level analysis applies statistical hypothesis
tests to the full training and evaluation distributions, and a
per-sample classification assigns each evaluation image
to one of four drift categories. The statistical approaches
utilised are directly deployed from the literature and presented in the following
subsections.



\subsubsection{Population-Level Statistical Tests}
\label{sec:pop_tests}

Four hypothesis tests are applied to characterise the nature
of the overall distributional shift between the training and
test/fine-tuning sets.

\textbf{MMD-RBF.} We utilize the Maximum Mean Discrepancy with an 
RBF kernel (MMD-RBF)\cite{gretton-borgwardt:2012} serves as a non-parametric 
two-sample test on the latent distributions.
The squared MMD is computed
between the training's set and fine-tuning's set latent vectors for each
model, with significance assessed via a permutation test with
$n_{\mathrm{perm}} = 300$ permutations. A significant result
indicates that the latent distributions are non-identical,
confirming the presence of data-level shift.

\textbf{Confidence Drift Test.} A Kolmogorov--Smirnov (KS) test\cite{kolmogorov:1933,smirnov:1948}
is applied to the prediction confidence distributions of the training and
fine-tuning sets for each model. A significant result (one
where the null hypothesis that the two distributions are drawn
from the same underlying distribution is rejected at $p < 0.05$)
indicates that the model's output certainty has shifted,
consistent with concept drift.

\textbf{HAS Margin Drift Test.} Building directly on the
angular margin signal defined in Equation~\ref{eq:margin},
a two-sample KS test is applied to the margin distributions
of the training and fine-tuning sets,assesing whether the population-level 
boundary proximity has shifted. Since the HASeparator penalty explicitly
maximises $\Delta_i$ for all training samples, the training
margin distribution represents the geometric baseline of a
well-adapted model. A significant drop in the fine-tuning
set margin ($\bar{\Delta}^{\mathrm{train}} - \bar{\Delta}^{\mathrm{ft}} > 0$)
indicates that embeddings in the fine-tuning set lie, on
average, closer to decision boundaries than training
embeddings.
That is a population-level geometric signature of
concept drift that is not accessible to the baseline model,
which operates in an unconstrained Euclidean space without
a meaningful notion of boundary proximity.

\textbf{HAS Boundary Direction Test.} Building on the
closest-boundary class signal defined in
Equation~\ref{sec:closest_boundary}, a $\chi^2$
goodness-of-fit test is applied to the distribution of
the closest-boundaryclass indices, comparing
the fine-tuning set distribution against the training set
reference. Whereas the HAS Margin Drift Test quantifies
\emph{how much} the population has drifted \emph{toward which boundaries}.


\subsubsection{Per-Sample Drift Classification}
\label{sec:persample}

Each fine-tuning sample is independently assigned a drift
category based on model-specific thresholds based on the
training statistics from Section~\ref{sec:extraction}.

Thus, each fine-tuning image is assigned to one of four mutually exclusive categories
based on the joint outcome of the data-drift and concept-drift
indicators: \emph{In-Distribution}, \emph{Pure Data Drift}, \emph{Pure Concept Drift},  and \emph{Full Drift}.

\textbf{Baseline.} A sample is classified as data-drifted if
its Euclidean distance from the global centroid exceeds
$\tau_d = \mu_d + k\,\sigma_d$, and as concept-drifted if
its confidence falls below
$\tau_q = \max(\mu_q - k\,\sigma_q,\ 0)$,
where $k = 2.0$.

\textbf{HAS Model.} Data drift is detected via the
class-conditional cosine distance threshold of
Equation~\ref{eq:cos_dist}, and concept drift via the
angular margin threshold of Equation~\ref{eq:margin_thresh},
both with $k = 2.0$.
Both thresholds follow the same $k\sigma$ criterion applied to
the training distribution of each signal, but in opposite
directions.
Data drift is detected when $d_\mu(i)$ exceeds
the upper tail of its training distribution, indicating the
embedding has moved away from its class centroid, while concept drift
is detected when $\Delta_i$ falls below the lower tail,
indicating the embedding has moved toward a decision boundary.

\subsubsection{Drift Direction and Hierarchical Analysis}
\label{sec:direction}

For concept-drifted samples, the closest-boundary class
$\hat{y}_i^{(2)}$ is recorded as the predicted drift
destination. A cross-tabulation of true class versus predicted
drift destination reveals the dominant pairwise confusion
direction for each class under distribution shift.
Complementarily, HAS drift direction matrices report the
normalised flow of samples from each true class to each
closest-boundary class, computed separately for the training
set, the full fine-tuning set, and the concept-drifted subset.
A shift in the direction distribution between training and
fine-tuning sets is a hallmark of concept drift that cannot
be captured by scalar confidence measures alone.

Both models are additionally evaluated at the sub-category
level using KS tests on the latent distance and confidence
distributions per sub-category. The fraction of drifted
sub-categories within each class provides a class-level drift
verdict and allows comparisons between models regarding which
class distributions are most affected by the domain gap.




